\scanpage{31}

$\widehat A'$ en un schéma formel abélien $\widehat A$.

Le th. de comparaison alg. formel nous ramène à prouver que $A$ est algébrique.

On est ainsi ramené à prouver le

\uline{Th.} Soit $\deleted{\unsure{B}}R$ anneau loc noeth complet, $S=\operatorname{Spec} R$, $\widehat A$ un schéma \inserted{abélien formel} abélien sur $\widehat S=\operatorname{Spf}\unsure{R|B}$. Conditions équivalentes

(i) $\widehat A$ algébri\deleted{\unsure{fiable|sable}}que, i.e. $\exists$ schéma abélien $A$ sur $S$, tel que $\widehat A$ soit isom. à son complété formel \unsure{le long} de $A_s$.

(ii) $\exists$ morphisme propre \uline{surjectif} \deleted{\illegible{1}} $f:S'\to S$, \deleted{\illegible{1}} tel que l'image inverse $\widehat A'$ de $\widehat A$ sur $\widehat S'$ soit algébrique.

(iii) Pour toute \deleted{algèbre finie $B$ sur $A$} quotient \deleted{\unsure{fini}} \inserted{intègre} $B$ de $A$, de dim $1$, si $B'$ est le normalisé de $B$, et $\deleted{\unsure{\widehat B}}$ $\widehat T'=\operatorname{Specf} B'$, \unsure{---} $\widehat A_{\widehat T'}$ est algébrique.

% slanted
\margin{\uline{NB} La démonstration montre qu'on peut se borner \uline{aux} $x$, $x\in S$, $\dim\{x\}=1$, qui ne sont pas contenus dans \deleted{\unsure{tous}} $\bigcup_i T_i$, où les $T_i$ sont \deleted{\unsure{des}} fermés dans $S$ \unsure{et leur|de leur} famille \uline{dénombrable}.}

\deleted{N.B. On peut se borner à prouver le th. si $\dim R=1$, et on \unsure{utilise} \unsure{alors|là} le \unsure{théorème} de \unsure{descente} \unsure{d'effectivité|l'effectivité} de descente fait \unsure{précédemment}.}

\uline{Dém} \quad (i) $\Rightarrow$ (ii) est trivial.

(ii) $\Rightarrow$ (iii) \quad En introduisant $T'=S'\times_S T$, on voit que $\widehat T'\unsure{-|\to}\widehat T$ ``rend $\widehat A_{\widehat T}$ algébrique''. Or on peut trouver, comme $T$ est un trait complet, un \deleted{\unsure{un|autre}} trait $T_1$ fini \deleted{et plat} sur $T$, et un $T$-morphisme $T_1-T'$. Donc $\widehat A_{\widehat T_1}$ est algébrique, donc définit un schéma abélien $A_{T_1}$ sur $T_1$. Ce dernier est muni d'une donnée de descente rel à $T_1-T$ [grâce au th. de comparaison], il suffit \deleted{\unsure{alors}} de voir que la donnée est effective, ce qui est évident, car $T_1-T$ est \inserted{fid} plat fini, et $A_{1}$ est projectif.
% slanted
\margin{\deleted{Pour utiliser \unsure{le} th. de \unsure{N.O.|N.O.S.}, \unsure{la série|le point} \unsure{singulière|singulier} du lieu singulier est \unsure{fermé|fourni}, \unsure{ceci|ce} est \unsure{évidemment} à \unsure{regarder}.}}
% archive sheet number: 30
